\subsection{器官移植受者的性生活：等待、移植与新常态}

肾、肝、心、肺等实体器官移植受者的性健康，常被铺天盖地的随访议程淹没——"活下来、不排斥"当然优先，但出院后的生活质量里，亲密从来都该在清单上。前一节《慢性肾衰竭与性功能》已从肾病视角谈过肾移植对性功能的改善；本节面向所有移植受者，梳理从等待期到"新常态"的常见议题。

\vspace{0.5em}
\noindent\textbf{等待期：欲望普遍"让位于生存"}：

\begin{itemize}
	\item 终末期器官疾病带来的疲劳、贫血、疼痛与抑郁，使性欲与性功能普遍下降——透析患者的具体情况见前一节
	\item 等待供体的不确定感本身就是最强的性欲抑制剂；伴侣之间"一方不敢要求、一方不敢开口"十分常见——把这个话题在随访时直接摆上桌面，反而轻松
\end{itemize}

\vspace{0.5em}
\noindent\textbf{术后早期}：

\begin{itemize}
	\item 以伤口愈合与体力恢复为先：心脏移植者胸骨愈合一般需 6--8 周，恢复性生活的体能自评可参照心血管事件后的标准（能连爬两层楼或快步走而不胸痛、气促）；伤口、引流管与导尿管的保护见"带管时的注意"相关内容
	\item 免疫抑制强化期内（术后最初数月）感染风险最高——恢复性生活后\textbf{安全套始终必要}（理由见下文）
\end{itemize}

\vspace{0.5em}
\noindent\textbf{长期管理中的性相关议题}：

\begin{itemize}
	\item \textbf{多数受者的性功能较终末期改善}：移植成功纠正代谢与血流动力学紊乱后，性欲与性能力多会回升；个体差异大，与移植物功能及基础神经血管病变相关（肾移植部分见前一节）
	\item \textbf{药物留下的"外貌与自信账单"}：糖皮质激素的满月脸与体重增加、环孢素/他克莫司的多毛与牙龈增生、震颤——这些变化不直接作用于性器官，却直接作用于\textbf{身体意象}与自信；用药方案常有调整空间，值得在复诊时提出，而不是默默忍受
	\item \textbf{PDE5 抑制剂（西地那非等）与移植药的相互作用}：与钙调磷酸酶抑制剂同经 CYP3A4 代谢，联用前需移植团队评估（血压、肾功能与血药浓度监测）——\textbf{切勿自行网购"壮阳药"}：成分不明产品与免疫抑制状态叠加，风险成倍放大
	\item \textbf{感染防护是长期课题}：免疫抑制状态下，经性传播的感染（HIV、乙肝/丙肝、梅毒、CMV 等）后果比常人更重——固定伴侣同查、新伴侣防护、随访筛查不可省略
\end{itemize}

\vspace{0.5em}
\noindent\textbf{生育与避孕}：

\begin{itemize}
	\item \textbf{男性}：移植后生育力多可恢复；但部分药物（如霉酚酸酯类）有精子毒性，\textbf{有生育计划者应在复诊时主动提出}，必要时调整方案；备孕前的准备与评估参照男性卷生育部分
	\item \textbf{女性}：一般建议移植后 1--2 年、原发病与移植物功能稳定后再妊娠——妊娠对免疫与循环系统都是重大挑战，需移植团队与高危产科共同管理；等待期与治疗期的\textbf{避孕}同样重要（意外妊娠对移植物与胎儿都是风险）——以屏障法为基础，激素避孕的选择遵移植团队
\end{itemize}

\vspace{0.5em}
\noindent\textbf{心理与"新常态"}：排斥恐惧、"我在吃免疫抑制剂"的自我认知、伴侣的过度保护（"你是病人，别累着"）都可能让亲密变形——移植受者的抑郁与焦虑筛查应包含性健康议题；移植专科随访联合心理科支持，是最经济的路径。

\begin{tcolorbox}[colback=pink!5!white,colframe=red!50!black,title=贴心小叮咛]
给随访清单加一行"性 / 亲密"——移植团队见过大量同类问题，这是他们知识范围内的正常话题。你不开口，它就永远排在议程最后；你开了口，多数困扰都有办法。
\end{tcolorbox}
